
Language models trained for code generation are typically evaluated on single metrics such as correctness (pass@k on benchmarks like HumanEval). Different alignment methods---execution-based training using pass/fail signals, preference-based training using human judgments, or no alignment---produce models with varying performance profiles. Recent work demonstrates that explicit multi-objective training can optimize for multiple dimensions simultaneously, but the implicit effects of standard alignment methods across quality dimensions remain unmeasured.

This work proposes a diagnostic framework for detecting alignment method signatures through multi-dimensional performance profiling. If alignment methods act as implicit objective functions---with execution feedback optimizing for correctness and preference feedback potentially optimizing differently---then models should exhibit detectable patterns in multi-dimensional performance space. Measuring models across correctness, complexity, and efficiency dimensions and applying clustering analysis could reveal these signatures.

\subsubsection{Scope and Contributions}

This paper presents a proof-of-concept validation of the diagnostic methodology. The primary contribution is demonstrating that signature detection is technically feasible through:

\begin{enumerate}
\item A diagnostic framework measuring models across three dimensions (correctness, complexity, efficiency) with PCA-based clustering and Cohen's d effect size analysis
\item Proof-of-concept validation using simulated data demonstrating the framework detects patterns when present (simulated Cohen's d=7.835)  
\item Minimal real-model validation (3 models, 10 HumanEval+ tasks, 300 generations) confirming the framework can profile actual models
\item Identification of critical limitations requiring resolution: model scale confounds, small sample size, and need for full-scale real-model validation
\end{enumerate}

The work does not claim that alignment methods create signatures in practice---only that if signatures exist, this methodology could detect them. Real-model validation at scale is required to test whether actual alignment methods produce detectable signatures.

\subsubsection{The Problem}

Code generation models aligned with different methods (SelfCodeAlign achieving 67.1\% on HumanEval+, preference-based methods, unaligned baselines) show varying performance. These differences may reflect systematic ''signatures''---models optimizing for whatever their training feedback measures. However, no prior work systematically measures alignment method effects across multiple quality dimensions via reverse-engineering (inferring objectives from outputs).

\subsubsection{The Gap}

Measuring signatures requires multi-dimensional evaluation infrastructure, cross-method comparison, and statistical frameworks for pattern detection. Most work optimizes known objectives forward (design multi-objective training) rather than inferring unknown objectives backward (detect implicit biases from outputs). Without signature detection, practitioners cannot diagnose model optimization biases or select methods based on quality profiles.

\subsubsection{The Insight}

Alignment methods may leave detectable signatures. If execution feedback optimizes for correctness, execution-trained models should occupy distinct regions in 3D space defined by (correctness, complexity, efficiency). Multi-dimensional evaluation combined with PCA and Cohen's d effect size could enable signature detection.

The remainder of this paper is organized as follows. Section 2 reviews related work. Section 3 describes the methodology. Section 4 presents experimental design. Section 5 reports proof-of-concept results. Section 6 discusses limitations and future work. Section 7 concludes.
